\begin{name}
	{\tenchude}
	{TOÁN 10}
	{LỚP TOÁN THẦY PHÁT}
	{Cuộc sống như chiếc gương, bạn chỉ nhận được kết quả tốt đẹp nhất khi mỉm cười với nó.}
\end{name}
\Opensolutionfile{ansbook}[ans/ansbookDe1]
\TN
\Opensolutionfile{ans}[ans/ansDe1-TN1]
\begin{ex}%[0-HK1-CT-1-BuiThiXuan-TPHCM-2324]%[VN-MT-6, VM013]%[0D6N3-3]
Mẫu số liệu sau thống kê số xe đạp đã bán được trong năm $2022$ của cửa hàng $A$ như sau:
\begin{center}
\begin{tabular}{cccccccccccc}
$10$&$7$&$8$&$3$&$7$&$15$&$25$&$16$&$17$&$9$&$8$&$7$\\
\end{tabular}
\end{center}
Hãy xác định trung vị của mẫu số liệu trên.
\choice
{$7$}
{$8$}
{$20$}
{\True $8{,}5$}
\loigiai{
Sắp xếp mẫu số liệu thành dãy không giảm, ta được
\begin{center}
\begin{tabular}{cccccccccccc}
$3$&$7$&$7$&$7$&$8$&$8$&$9$&$10$&$15$&$16$&$17$&$25$\\
\end{tabular}
\end{center}
Do cỡ mẫu của mẫu số liệu trên bằng $12$ nên trung vị của mẫu số liệu là \[\dfrac{8+9}{2}=8{,}5.\]
}
\end{ex}

\begin{ex}%[HK1, THPT Lê Quảng Chí, 2023]%[TVN-001, 10-11EX-HK1-2324]%[0D6N4-2]
Kết quả dự báo nhiệt độ cao nhất trong 10 ngày cuối tháng 12 năm 2022 ở Yên Bái thu được như sau
\[
\begin{array}{lllllllllll}
20 & 21 & 20 & 19 & 22 & 17 & 18 & 18 & 14 & 16 & \left({}^{\circ}\mathrm{C}\right) .
\end{array}
\]
Khoảng biến thiên $R$ của mẫu số liệu trên là
\choice
{$R=7$}
{\True $R=8$}
{$R=6$}
{$R=5$}
\loigiai{Giá trị lớn nhất trong mẫu số liệu là $22$.\\
Giá trị nhỏ nhất trong mẫu số liệu là $14$.\\
Khoảng biến thiên là $R=22-14=8$.}
\end{ex}

\begin{ex}%[10-11-EX-GK2-2425]%[Nguyễn Chiến]%[0D8H1-3]
Một mạng đường giao thông nối các tỉnh A, B, C, D, E, F và G như hình vẽ, trong đó số viết trên mỗi cạnh nối $2$ tỉnh với nhau thể hiện số con đường nối $2$ tỉnh đó. Hỏi có bao nhiêu cách đi từ tỉnh A đến tỉnh G?
\begin{center}
\begin{tikzpicture}[scale=0.8, line join = round, line cap = round]
\tikzset{label style/.style={font=\footnotesize}}
\def \canhx{3}
\def \canhy{1.5}
\coordinate[label= left: $A$] (A) at (0,0);
\coordinate[label= above: $B$] (B) at (\canhx,\canhy);
\coordinate[label= below: $C$] (C) at (\canhx,-\canhy);
\coordinate[label= above: $D$] (D) at (2*\canhx,0);
\coordinate[label= above: $E$] (E) at (3*\canhx,\canhy);
\coordinate[label= below: $F$] (F) at (3*\canhx,-\canhy);
\coordinate[label= right: $G$] (G) at (4*\canhx,0);
\draw[] ($(A)!0.5!(B)$) node[above]  {$2$} ($(B)!0.5!(D)$) node[above]  {$3$} ($(D)!0.5!(E)$) node[above]  {$5$} ($(E)!0.5!(G)$) node[above]  {$7$} ($(A)!0.5!(C)$) node[below]  {$8$} ($(C)!0.5!(D)$) node[below]  {$6$} ($(D)!0.5!(F)$) node[below]  {$3$} ($(F)!0.5!(G)$) node[below]  {$4$};
\draw[] (A)--(B)--(F)--(G)--(E)--(C)--(A);
\foreach \i in {A,B,C,D,E,F,G} \draw[fill=black] (\i) circle(1.2pt)  ; % vẽ điểm
\end{tikzpicture}
\end{center}
\choice
{$361$}
{\True $2\, 538$}
{$262\,080$}
{$786$}
\loigiai{
Ta xét các trường hợp sau
\begin{itemize}
\item Trường hợp 1: Đi theo đường A-B-D-E-G.\\
Số cách đi là $2 \cdot 3 \cdot 5 \cdot 7=210$ cách.
\item Trường hợp 2: Đi theo đường A-B-D-F-G.\\
Số cách đi là $2 \cdot 3 \cdot 3 \cdot 4=72$ cách.
\item Trường hợp 3: Đi theo đường A-C-D-E-G.\\
Số cách đi là $8 \cdot 6 \cdot 5 \cdot 7=1 \, 680$ cách.
\item Trường hợp 4: Đi theo đường A-C-D-F-G.\\
Số cách đi là $8 \cdot 6 \cdot 3 \cdot 4=576$ cách.
\end{itemize}
Vậy có $210+72+1 \, 680+676=2 \, 538$ cách đi từ tỉnh A đến tỉnh G.
}
\end{ex}

\begin{ex}%[0D8H2-3]
Có bao nhiêu số tự nhiên có sáu chữ số khác nhau từng đôi một, trong đó chữ số $ 5$ đứng liền giữa hai chữ số $ 1$ và $ 4$?
\choice
{$249$}
{\True $1500$}
{$3204$}
{$2942$}
\loigiai{
Chữ số $ 5$ đứng liền giữa hai chữ số $ 1$ và $ 4$ nên ta có thể có $ 154$ hoặc $ 451$.\\
Gọi số cần tìm là $\overline{abc}$, sau đó ta chèn thêm $ 154$ hoặc $451$ để có được số gồm 6 chữ số cần tìm.
\begin{itemize}
\item \textbf{Trường hợp 1:} $ a\ne 0$, số cách chọn $ a$ là $ 6$, số cách chọn $ b$ và $ c$ là $ \mathrm{A}_6^2$, sau đó chèn $ 154$ hoặc $ 451$ vào $ 4$ vị trí còn lại nên có $ 6\cdot\mathrm{A}_6^2\cdot4\cdot2$ cách.
\item \textbf{Trường hợp 2:} $ a=0$, số cách chọn $ a$ là 1, số cách chọn $ b$ và $ c$ là $ \mathrm{A}_6^2$, sau đó chèn $ 154$ hoặc $ 451$ vào vị trí trước $ a$ có duy nhất 1 cách nên có $ \mathrm{A}_6^2\cdot2$ cách.
\end{itemize}
Vậy có $ 6\cdot\mathrm{A}_6^2\cdot4\cdot2+
\mathrm{A}_6^2\cdot2=1500$.}
\end{ex}

\begin{ex}%[0D8H2-4]
Một nhóm gồm $10$ học sinh, trong đó có $4$ học sinh nam và $6$ học sinh nữ. Hỏi có bao nhiêu cách chọn ra $2$ học sinh từ nhóm đó sao cho $2$ học sinh được chọn có cả nam và nữ?
\choice
{$20$}
{$36$}
{$10$}
{\True $24$}
\loigiai{
Để chọn được $2$ học sinh được chọn có cả nam và nữ ta chọn $1$ nam và $1$ nữ.\\
Số cách chọn thoả mãn yêu cầu đề bài là $4\cdot 6=24$.
}
\end{ex}

\begin{ex}%[0D8H2-6]
Cho $10$ điểm phân biệt và không có ba điểm nào thẳng hàng. Hỏi có bao nhiêu véc-tơ khác $\overrightarrow{0}$ có điểm đầu và điểm cuối là hai trong số $10$ điểm đã cho?
\choice
{\True $90$}
{$45$}
{$10$}
{$20$}
\loigiai{
Số véc-tơ khác $\overrightarrow{0}$ có điểm đầu và điểm cuối là hai trong số $10$ điểm đã cho là $\mathrm{A}_{10}^2=90$.
}
\end{ex}

\begin{ex}%[0D8H3-2]
Khai triển $(x+1)^4$ ta được
\choice
{\True $x^4+4x^3+6x^2+4x+1$}
{$x^4-4x^3+6x^2-4x+1$}
{$x^4+6x^3+8x^2+6x+1$}
{$x^4+4x^3+8x^2+4x+1$}
\loigiai{
Ta có $(x+1)^4=\mathrm{C}_4^0\,x^4+\mathrm{C}_4^1\,x^3+\mathrm{C}_4^2\,x^2+\mathrm{C}_4^3\,x+\mathrm{C}_4^4=x^4+4x^3+6x^2+4x+1$.
}
\end{ex}

\begin{ex}%[0-HK1-KN-1-PhanDinhPhung-HaNoi-2324]%[VN-MT-6, VM019]%[0H9H1-1]
Trong mặt phẳng tọa độ $Oxy$, cho $\overrightarrow{OM}=\overrightarrow{i}+2\overrightarrow{j}$, $\overrightarrow{ON}=2\overrightarrow{i}$. Tìm tọa độ vectơ $\overrightarrow{MN}$.
\choice
{\True $\overrightarrow{MN}=(1;-2)$}
{$\overrightarrow{MN}=(1;2)$}
{$\overrightarrow{MN}=(-1;-2)$}
{$\overrightarrow{MN}=(-1;2)$}
\loigiai{
Tọa độ điểm $M(1;2)$ và $N(2;0)$ nên $\overrightarrow{MN}=(1;-2)$.
}
\end{ex}

\begin{ex}%[0H9H3-2]
Trong mặt phẳng $Oxy$, cho điểm $M\left(1;2\right)$. Gọi $A$, $B$ là hình chiếu của $M$ lên $Ox$, $Oy$. Viết phương trình đường thẳng $AB$.
\choice
{$x+2y-1=0$}
{$2x+y+2=0$}
{\True $2x+y-2=0$}
{$x+y-3=0$}
\loigiai{
Ta có hình chiếu của điểm $M\left(1;2\right)$ lên $Ox$, $Oy$ lần lượt là $A$ và $B$.\\
Do đó phương trình đường thẳng $AB$ là $\dfrac{x}{1}+\dfrac{y}{2}=1\Leftrightarrow 2x+y-2=0$.
}
\end{ex}

\begin{ex}%[0H9H4-2]%[Dự án đề kiểm tra Toán 10 HKII NH23-24- Trung Anh]%[THPT Việt Đức - Hà Nội]
Trong mặt phẳng với hệ trục tọa độ $Oxy$, đường tròn $(C)$ có tâm $I(1;-2)$ và tiếp xúc với trục $Ox$ có phương trình là
\choice
{$(C): (x+1)^2+(y-2)^2=1$}
{\True $(C): (x-1)^2+(y+2)^2=4$}
{$(C): (x-1)^2+(y+2)^2=1$}
{$(C): (x+1)^2+(y-2)^2=4$}
\loigiai{
Đường tròn $(C)$ có tâm $I(1;-2)$ và tiếp xúc với trục $Ox$ nên suy ra \[R=|-2|=2.\]
Đường tròn $(C)$ có phương trình là \[(x-1)^2+(y+2)^2=4.\]
}
\end{ex}

\begin{ex}%[soạn để GHK, HK]%[Huỳnh Xuân Tín]%[0H9H5-2]
Cho elip $(E)$ có một đỉnh $A_2(7;0)$, một tiêu điểm $F_2(5;0)$.  Lập phương trình chính tắc của $(E)$.
\choice
{\True $\dfrac{x^2}{49}+\dfrac{y^2}{24}=1$}
{$\dfrac{x^2}{24}+\dfrac{y^2}{49}=1$}
{$\dfrac{x^2}{49}+\dfrac{y^2}{25}=1$}
{$\dfrac{x^2}{49}-\dfrac{y^2}{24}=1$}
\loigiai{
Ta có đỉnh $A_2(7;0)$, một tiêu điểm $F_2(5;0)$, suy ra $a=7$,  $c=5$ và $b^2=a^2-c^2=24$, do đó  $\dfrac{x^2}{49}+\dfrac{y^2}{24}=1$.
}
\end{ex}

\begin{ex}%[0H9H5-4]
Tổng các giá trị $m$ để hypebol $(H)\colon \dfrac{x^2}{4}-\dfrac{y^2}{m^2-2m-7}=1$ có tiêu điểm trung với tiêu điểm của elip $(E)\colon\dfrac{x^2}{16}+\dfrac{y^2}{9}=1$ là
\choice
{$4$}
{\True $2$}
{$5$}
{$6$}
\loigiai{
Tiêu điểm của elip $(E)\colon \dfrac{x^2}{16}+\dfrac{y^2}{9}=1$ là
\[F_{1;2}=\left(\pm\sqrt{16-9};0\right)=\left(\pm\sqrt{7};0\right).\]
Tiêu điểm của hypebol $(H)\colon \dfrac{x^2}{4}-\dfrac{y^2}{m^2-2m-7}=1$ là \[F'_{1;2}=\left(\pm\sqrt{a^2+b^2};0\right)=\left(\pm\sqrt{4+m^2-2m-7};0\right).\]
\begin{eqnarray*}
&\Rightarrow& \left(\pm\sqrt{4+m^2-2m-7};0\right)=\left(\pm\sqrt{7};0\right)\\
&\Rightarrow& 4+m^2-2m-7=7\\
&\Rightarrow& m^2-2m-10=0\\
&\Rightarrow& \hoac{&m=1+\sqrt{11}\\ &m=1-\sqrt{11}.}
\end{eqnarray*}
Vậy tổng của giá trị $m$ để thỏa yêu cầu bài toán là $2$.
}
\end{ex}
\TNTF
\begin{ex}%[0D8H2-4]
Một lớp có $24$ học sinh nam và $16$ học sinh nữ. Chọn ra $3$ học sinh làm ban cán sự lớp.
\choiceTF
{Có $59280$ cách chọn tùy ý}
{\True Có $4416$ cách chọn sao cho trong đó có $2$ học sinh nam}
{\True Có $9320$ cách chọn sao cho trong đó có ít nhất $1$ học sinh nam}
{\True Có $6440$ cách chọn sao cho trong đó có nhiều nhất $1$ học sinh nữ}
\loigiai{
\begin{itemchoice}
\itemch Sai vì\\
Mỗi cách chọn $3$ học sinh trong $40$ học sinh là một tổ hợp chập $3$ của $40$. Số cách chọn $3$ học sinh làm ban cán sự của lớp là $\mathrm{C}_{40}^3=9880$.
\itemch Đúng vì\\
Số cách chọn $2$ học sinh nam trong $24$ học sinh nam là $\mathrm{C}_{24}^2=276$. \\
Số cách chọn $1$ học sinh nữ trong $16$ học sinh nữ là $\mathrm{C}_{16}^1=16$. \\
Vậy số cách chọn $3$ học sinh làm ban cán sự lớp sao cho trong đó có $2$ học sinh nam là ${276\cdot 16=4416}$.
\itemch Đúng vì\\
Số cách chọn $3$ học sinh làm ban cán sự của lớp là $\mathrm{C}_{40}^3=9880$. \\
Số cách chọn $3$ học sinh nữ làm ban cán sự của lớp là $\mathrm{C}_{16}^3=560$. \\
Vậy số cách chọn $3$ học sinh làm ban cán sự của lớp sao cho trong đó có ít nhất $1$ học sinh nam là
$9880-560=9320$.
\itemch Đúng vì\\
Số cách chọn $3$ học sinh nam là $\mathrm{C}_{24}^3=2024$. \\
Số cách chọn $1$ học sinh nữ và $2$ học sinh nam: $\mathrm{C}_{16}^1\cdot\mathrm{C}_{24}^2=4416$. \\
Vậy số cách chọn $3$ học sinh làm ban cán sự của lớp sao cho trong đó có ít nhất $1$ học sinh nam là
$2024+4416=6440$.
\end{itemchoice}
}
\end{ex}

\begin{ex}%[Dự án 18 - TeamTeXHoa - Phạm Quốc Toàn]%[0H9H4-5]
Cho đường thẳng $d \colon x-my-2=0$ và đường tròn $\left( C \right) \colon \left( x-3 \right)^{2}+\left( y-1 \right)^{2}=4$. Các khẳng định sau đúng hay sai?
\choiceTF
{Với $m=2$ thì $d$ và đường tròn $\left( C \right)$ tiếp xúc}
{\True Với $m=1$ thì $d$ đi qua tâm đường tròn $\left( C \right)$}
{Khoảng cách từ tâm đường tròn $\left( C \right)$  đến đường thẳng $d$ luôn lớn hơn $\sqrt{2}$}
{\True Đường thẳng $d$ luôn cắt đường tròn $\left( C \right)$}
\loigiai{
$\left( C \right) \colon x^2+y^2-2x+4y-11=0$ có tâm $I\left( 3;1 \right)$ bán kính $R=2$.
\begin{itemchoice}
\itemch Với $m=2$ thì phương trình đường thẳng $d$ trở thành $x-2y-2=0$. \\
Khi đó $d\left( I,\Delta  \right)=\dfrac{\left| 3-2\left( -1 \right)-2 \right|}{\sqrt{1^2+2^2}}=\dfrac{3}{\sqrt{5}}<R$ . Do đó $d$ cắt $\left( C \right)$ chứ không tiếp xúc với  $\left( C \right)$.
\itemch Với $m=1$ thì phương trình đường thẳng $d$ trở thành $x-y-2=0$. \\
Ta có $3-1-2=0$ nên điểm $I\left( 3;1 \right)$ thuộc đường thẳng $d$.
\itemch Ta có $d\left( I,\Delta  \right)=\dfrac{\left| 3-m\left( -1 \right)-2 \right|}{\sqrt{1^2+m^2}}=\dfrac{\left| m+1 \right|}{\sqrt{m^2+1}}$. \\
Áp dụng bất đắng thức, ta có $\left( m+1 \right)^{2} \le 2\left( m^2+1 \right)$, $\forall m \Rightarrow \dfrac{\left| m+1 \right|}{\sqrt{m^2+1}} \le \sqrt{2}$, $\forall m$. \\
Do đó $d\left( I,\Delta  \right)\le \sqrt{2}$, $\forall m$.
\itemch $d\left( I,\Delta  \right) \le \sqrt{2}<2=R$, $\forall m$. Vậy đường thẳng $d$ luôn cắt $\left( C \right)$.
\end{itemchoice}
}
\end{ex}
\TNSA
\begin{ex}%[0D8H3-4]
Tính tổng các hệ số của số hạng chứa $x^n$ với $n\in\mathbb{N}$ trong khai triển nhị thức Newton $\left(x^2+\dfrac{2}{x}\right)^5$, với $x\ne0$.
\shortans[]{$131$}
\loigiai{Ta có
\begin{eqnarray*}
\left(x^2+\dfrac{2}{x}\right)^5&=&\mathrm{C}_5^0\left(x^2\right)^5+\mathrm{C}_5^1\left(x^2\right)^4\left(\dfrac{2}{x}\right)+\mathrm{C}_5^2\left(x^2\right)^3\left(\dfrac{2}{x}\right)^2+\mathrm{C}_5^3\left(x^2\right)^2\left(\dfrac{2}{x}\right)^3+\mathrm{C}_5^4\left(x^2\right)\left(\dfrac{2}{x}\right)^4+\mathrm{C}_5^5\left(\dfrac{2}{x}\right)^5\\
&=&x^{10}+10x^7+40x^4+80x+\dfrac{80}{x^2}+\dfrac{32}{x^5}.
\end{eqnarray*}
Vậy tổng các hệ số cần tìm là $T=1+10+40+80=131$.}
\end{ex}

\begin{ex}%[0D8V2-4]
Một đội thanh niên tình nguyện có $15$ người gồm $12$ nam và $3$ nữ. Hỏi có bao nhiêu cách phân công đội thanh niên tình nguyện đó về giúp đỡ $3$ tỉnh miền núi, sao cho mỗi tỉnh có $4$ nam và $1$ nữ?
\shortans{60}
\loigiai{
Ta có $\mathrm{C}_3^1\cdot\mathrm{C}_{12}^4$ cách phân công các thanh niên về tỉnh thứ nhất.\\
Với mỗi cách này thì có $\mathrm{C}_2^1 \cdot \mathrm{C}_8^4$ cách phân công số thanh niên còn lại về tỉnh thứ hai.\\
Với mỗi cách phân công trên thì có $\mathrm{C}_1^1\cdot\mathrm{C}_{4}^4$ cách phân công số thanh niên còn lại về tỉnh thứ ba.\\
Do đó, ta có $\mathrm{C}_3^1\cdot\mathrm{C}_{12}^4\cdot\mathrm{C}_2^1\cdot\mathrm{C}_{8}^4\cdot\mathrm{C}_1^1\cdot\mathrm{C}_{4}^4=207900$ cách phân công thỏa mãn đề bài.\\
Gọi số cần lập có dạng $abc$.\\
Chọn $c$ với $c \in\{4;6;8\}$ có $3$ cách.\\
Chọn $\overline{ab}$ có $\mathrm{A}_5^2$ cách.\\
Theo quy tắc nhân, ta có $3 \cdot \mathrm{A}_5^2=60$ số tự nhiên thỏa mãn.
}
\end{ex}

\begin{ex}%[0H9V4-1]
Trong mặt phẳng toạ độ $Oxy$, đường tròn $(C)$ đi qua hai điểm $A(-2;3)$, $B(3;0)$ và có tâm thuộc đường thẳng $d: x-2=0$. Tính chu vi đường tròn $(C)$ \textit{(kết quả làm tròn đến hàng phần trăm).}
\shortans{4,12}
\loigiai{
Vì đường tròn $(C)$ có tâm thuộc đường thẳng $d: x-2=0$ nên toạ độ tâm $I$ của đường tròn $(C)$ có dạng $(2;m)$, với $m\in\mathbb{R}$.

Mặt khác, đường tròn $(C)$ đi qua hai điểm $A(-2;3)$, $B(3;0)$ nên $IA=IB$.
\begin{eqnarray*} %Dễ căn hơn á
&\Leftrightarrow& IA^2=IB^2\\
&\Leftrightarrow& [2-(-2)]^2+(m-3)^2=(2-3)^2+(m-0)^2\\
&\Leftrightarrow& m^2-6m+25=m^2+1\\
&\Leftrightarrow &-6m+24=0\\
&\Leftrightarrow& m=4\\
&\Rightarrow& I(2;4)\Rightarrow R=IA=\sqrt{17}.
\end{eqnarray*}

Do đó, chu vi của đường tròn $(C)$ là $2\pi R=2\sqrt{17}\pi\approx 4{,}12$.
}
\end{ex}

\begin{ex}%[0H9V5-0]%[tex hóa đề ck2 - form 2025 - đợt 2- Duong Xuan Loi]
\immini{
Một nhà vòm chứa máy bay có mặt cắt hình nửa elip cao $5$ m, rộng $20$ m. Khoảng cách theo phương thẳng đứng từ một điểm cách chân tường $5$ m lên đến nóc nhà vòm bằng $\dfrac{a\sqrt{b}}{c}$ với $a$, $b$, $c$ là các số nguyên dương. Tính giá trị biểu thức $T=a+2b-c$.
}{
\begin{tikzpicture}[>=stealth,line join=round,line cap=round,font=\footnotesize,scale=.45]
%mái vòm
\draw[green!50!black] (8,0) arc(0:180:8 and 7.8);
\foreach \i in {1,2,3,4,5,6,7,8,9,10,11}
\draw[green!50!black] (0,0)--({15*(\i)}:8 and 7.8);
\fill[draw=green!50!black,white] (7.6,0)arc(0:180:7.6 and 7.5)--cycle;
\draw[green!50!black] (7.6,0) arc(0:180:7.6 and 7.5);

\draw[green!50!black] (6.7,0.8) arc(0:180:6.7 and 6.3);
\foreach \i in {1,2,3,4,5,6,7,8,9,10,11}
\draw[green!50!black] (0,0.8)--($(0,0.8)+({15*(\i)}:6.7 and 6.3)$);
\fill[draw=green!50!black,white] (6.4,0.8)arc(0:180:6.4 and 6.1)--cycle;
\draw[green!50!black] (6.4,0.8) arc(0:180:6.4 and 6.1);

\draw[green!50!black] (5.4,1.6) arc(0:180:5.4 and 4.8);
\foreach \i in {1,2,3,4,5,6,7,8,9,10,11}
\draw[green!50!black] (0,1.6)--($(0,1.6)+({15*(\i)}:5.4 and 4.8)$);
\fill[draw=green!50!black,white] (5.15,1.6)arc(0:180:5.15 and 4.6)--cycle;
\draw[green!50!black] (5.15,1.6) arc(0:180:5.15 and 4.6);

\foreach \i in {1,2,3,4,5,6,7,8}
\draw[green!50!black] (0,3.5)--({20*(\i)}:8 and 7.8);
\fill[draw=green!50!black,white] (4.1,2.4)arc(0:180:4.1 and 3.3)--cycle;
\draw[green!50!black] (4.1,2.4) arc(0:180:4.1 and 3.3);
\foreach \i in {1,2,3,4,5,6,7,8,9,10,11}
\draw[green!50!black] (0,2.4)--($(0,2.4)+({15*(\i)}:4.1 and 3.3)$);
\fill[draw=green!50!black,white] (3.85,2.4)arc(0:180:3.85 and 3.1)--cycle;
\draw[green!50!black] (3.85,2.4) arc(0:180:3.85 and 3.1);
\fill[middle color=cyan,opacity=0.05] (8,0) arc(0:180:8 and 7.8);
%cửa sau
\fill[draw=green!50!black,bottom color=cyan!20!white,top color=white] (3.85,2.4)arc(0:180:3.85 and 3.1)--cycle;
\draw[color=green!50!black] ($(0,2.4)+(60:3.85 and 3.1)$)|-(0,2.4) coordinate[pos=0.1](A);
\draw[color=green!50!black] ($(0,2.4)+(120:3.85 and 3.1)$)|-(0,2.4) coordinate[pos=0.1](B);
\draw[color=green!50!black] ($(0,2.4)+(60:3.85 and 3.1)$)|-($(0,2.4)+(30:3.85 and 3.1)$) ($(0,2.4)+(120:3.85 and 3.1)$)|-($(0,2.4)+(150:3.85 and 3.1)$) (A)--(B);
%sàn nhà
\fill[bottom color=gray!40!white,top color=white] (8,0)--(4.1,2.4)--(-4.1,2.4)--(-8,0)--cycle;
%bánh xe bên phải
\fill[color=black] (1,1.4)--(1.5,0.9)--++(0.1,0)--(1.1,1.4)--cycle;
\fill[color=black] (1.5,0.8)--++(0,0.7)--++(0.1,0)--++(0,-0.7)--++(0.3,0)--++(0,-0.5)--++(-0.2,0)--++(0,0.4)--++(-0.3,0)--++(0,-0.4)--++(-0.2,0)--++(0,0.5)--++(0.3,0)--cycle;
\fill[left color=gray] (2.6,1) circle(0.5 cm);
\fill[right color=black] (2.6,1) circle(0.4 cm);
\fill[left color=gray] (3.7,1.23) circle(0.4 cm);
\fill[right color=black] (3.7,1.23) circle(0.3 cm);
%bánh xe bên trái
\fill[xscale=-1,color=black] (1,1.4)--(1.5,0.9)--++(0.1,0)--(1.1,1.4)--cycle;
\fill[xscale=-1,color=black] (1.5,0.8)--++(0,0.7)--++(0.1,0)--++(0,-0.7)--++(0.3,0)--++(0,-0.5)--++(-0.2,0)--++(0,0.4)--++(-0.3,0)--++(0,-0.4)--++(-0.2,0)--++(0,0.5)--++(0.3,0)--cycle;
\fill[xscale=-1,left color=gray] (2.6,1) circle(0.5 cm);
\fill[xscale=-1,right color=black] (2.6,1) circle(0.4 cm);
\fill[xscale=-1,left color=gray] (3.7,1.23) circle(0.4 cm);
\fill[xscale=-1,right color=black] (3.7,1.23) circle(0.3 cm);
%cánh phải dưới
\fill[left color=gray,right color=cyan!50!black] (0.75,1.5)--(6,1.9)--++(0.5,0.35)--++(-0.2,0)--(6,2)--(0.7,1.9)--cycle;
\fill[left color=yellow!40!white, right color=cyan!50!black] (0.75,1.5)--(6,1.9)--($(0,1.3)+(-5:0.8 and 0.5)$)arc(0:-94:0.75 and 0.4)--($(0,1.5)+(-90:0.75 and 0.5)$) arc(-90:0:0.75 and 0.5)--cycle;
%cánh trái dưới
\fill[xscale=-1,right color=gray,left color=cyan!50!black] (0.75,1.5)--(6,1.9)--++(0.5,0.35)--++(-0.2,0)--(6,2)--(0.7,1.9)--cycle;
\fill[xscale=-1,right color=yellow!40!white, left color=cyan!50!black] (0.75,1.5)--(6,1.9)--($(0,1.3)+(-5:0.8 and 0.5)$)arc(0:-94:0.75 and 0.4)--($(0,1.5)+(-90:0.75 and 0.5)$) arc(-90:0:0.75 and 0.5)--cycle;
%cánh phải trên
\fill[left color=gray,right color=cyan!50!black] (0.7,2)--(3.1,2.15)--++(0.05,0.05)--($(0,2)+(15:0.7 and 0.5)$)--cycle;
%cánh trái trên
\fill[xscale=-1,right color=gray,left color=cyan!50!black] (0.7,2)--(3.1,2.15)--++(0.05,0.05)--($(0,2)+(15:0.7 and 0.5)$)--cycle;
%đuôi máy bay
\fill[bottom color=yellow!40!white,top color=cyan!40!black] ($(0,2)+(80:0.7 and 0.5)$)--(0,5)--($(0,2)+(100:0.7 and 0.5)$)--cycle;
%đầu máy bay
\fill[bottom color=black!70, top color=cyan] (0.7,2) arc(0:180:0.7 and 0.5)--(-0.75,1.5) arc(180:0:0.75 and 0.2)--cycle;
\fill[top color=black!70, bottom color=cyan!50!white](-0.75,1.5) arc(-180:0:0.75 and 0.5) arc(0:180:0.75 and 0.2)--cycle;
\fill[color=black] (0.75,1.5)--(0.7,2)--(0,1.9)--(-0.7,2)--(-0.75,1.5) arc(180:0:0.75 and 0.2)--cycle;
\fill[pattern color=gray, pattern=vertical lines] (0.75,1.5)--(0.7,2)--(0,1.9)--(-0.7,2)--(-0.75,1.5) arc(180:0:0.75 and 0.2)--cycle;
%bánh xe giữa
\fill[color=black] ($(0,1.3)+(-90:0.75 and 0.4)$)--(0.1,0.6)--++(0.1,0)--++(0,-0.3)--++(-0.1,0)--++(0,0.25)--($(0,1.25)+(-90:0.75 and 0.4)$)--cycle;
\fill[xscale=-1,color=black] ($(0,1.3)+(-90:0.75 and 0.4)$)--(0.1,0.6)--++(0.1,0)--++(0,-0.3)--++(-0.1,0)--++(0,0.25)--($(0,1.25)+(-90:0.75 and 0.4)$)--cycle;
\end{tikzpicture}
}
\shortans[0]{$9$}
\loigiai{
Chọn hệ trục tọa độ $Oxy$ với gốc tọa độ tại tâm đáy nhà vòm, trục tung thẳng đứng.
\begin{center}
\begin{tikzpicture}[scale=0.5, font=\footnotesize, line join=round, line cap=round,>=stealth]
\def\xmin{-14} \def\xmax{14}
\def\ymin{-7} \def\ymax{7}
\draw[->] (\xmin,0)--(\xmax,0) node [below]{$x$};
\draw[->] (0,\ymin)--(0,\ymax) node [left]{$y$};
\node at (0,0) [below left]{$O$};
\clip (\xmin+0.1,\ymin+0.1) rectangle (\xmax-0.1,\ymax-0.1);
\draw (0,0) ellipse (10cm and 5cm);
\fill (10,0) circle (2.0pt) node[above right]{$C(10;0)$}
(-10,0) circle (2.0pt) node[above left]{$B(-10;0)$}
(0,5) circle (2.0pt) node[above right]{$A(0;5)$};
\end{tikzpicture}
\end{center}
Nhà vòm có dạng nửa elip nên có phương trình chính tắc của elip là $\dfrac{x^2}{a^2}+\dfrac{y^2}{b^2}=1$ ($a, b>0$).\\
Ta có chiều cao của nhà vòm là $5$ m nên $OA=h=5$, chiều rộng của nhà vòm là $20$ m nên $BC=2OB=20$. Suy ra $OB=10$.\\
Ta có tọa độ các điểm: $C(10; 0)$ và $A(0; 5)$. Thay hai điểm này vào phương trình chính tắc, ta có\\
$\heva{&\dfrac{10^2}{a^2}+\dfrac{0^2}{b^2}=1\\&\dfrac{0^2}{a^2}+\dfrac{5^2}{b^2}=1}\Leftrightarrow\heva{&a=10\\&b=5\cdot.}$ \\
Suy ra phương trình miêu tả hình dáng nhà vòm là $\dfrac{x^2}{100}+\dfrac{y^2}{25}=1$.\\
Điểm cách chân tường $5$ m tương ứng cách tâm $5$ m (vì từ tâm vòm đến tường là $10$ m).\\
Thay $x=5$ vào phương trình $\dfrac{x^2}{100}+\dfrac{y^2}{25}=1$, ta tìm được $y=\dfrac{5\sqrt{3}}{2}$.\\
Khi đó: $\heva{&a=5\\&b=3\\&c=2}\Rightarrow T=a+2b-c=5+2\cdot 3-2=9$.}
\end{ex}

\TL


\begin{ex}%[0D8H3-2]
Sử dụng công thức nhị thức Newton khai triển biểu thức $\left(x-\sqrt{2}\right)^5$.
\loigiai{
Ta có
\begin{eqnarray*}
\left(x-\sqrt{2}\right)^5&=&\mathrm{C}_5^0x^5-\mathrm{C}_5^1x^4\left(\sqrt{2}\right)^1+\mathrm{C}_5^2x^3\left(\sqrt{2}\right)^2-\mathrm{C}_5^3x^2\left(\sqrt{2}\right)^3+\mathrm{C}_5^4x^1\left(\sqrt{2}\right)^4-\mathrm{C}_5^5\left(\sqrt{2}\right)^5 \\
&=&x^5-5\sqrt{2}x^4+20x^3-20\sqrt{2}x^2+20x-4\sqrt{2}.
\end{eqnarray*}
}
\end{ex}
\begin{ex}%[0D8H2-5]
Từ một hộp chứa $3$ quả cầu đỏ và $4$ quả cầu trắng. Có bao nhiêu cách lấy ngẫu nhiên $3$ quả cầu sao cho có ít nhất $1$ quả cầu đỏ?
% \shortans[]{$31$}
\loigiai{
Số cách lấy $3$ quả cầu bất kì từ hộp là $\mathrm{C}^3_7$.\\
Số cách lấy $3$ quả cầu đều có màu trắng là $\mathrm{C}^3_4$.\\
Số cách lấy ngẫu nhiên $3$ quả cầu sao cho có ít nhất $1$ quả cầu đỏ là $\mathrm{C}^3_7-\mathrm{C}^3_4=31$.
}
\end{ex}

\begin{ex}%[0H9V4-2]
Ba bạn $A$, $B$, $C$ ở trên sân chơi trò chơi như sau. Ba bạn đứng một chỗ, $A$ đi về hướng Tây $1$ m sau đó đi về hướng Bắc $4$ m, $B$ đi về phía đông $5$ m sau đó đi về phía bắc $6$ m, $C$ đi về phía đông $6$ m sau đó đi về phía bắc $3$ m. Đặt cờ ở vị trí nào để ba bạn cách đều cờ nhất?? Vị trí cờ tính từ vị trí ba bạn đứng cần đi theo hướng nào để đến được?
\loigiai{
Gắn hệ trục toạ độ $Oxy$ sao cho gốc toạ độ trùng với vị trí ba bạn đứng ban đầu. Khi đó
$A(-1 ; 4)$, $B(5 ; 6)$, $C(6 ; 3)$ là toạ độ tương ứng của $A$, $B$, $C$. Vị trí đặt cờ cách đều ba bạn nên là tâm đường tròn đi qua $A$, $B$, $C$.
Gọi $I(x ; y)$ là tâm đường tròn ngoại tiếp tam giác $ABC$. Ta có $IA=IB$ và $IB=IC$ nên
\begin{align*}
&~ \heva{& IA^2=IB^2\\&IB^2=IC^2}\\
\Leftrightarrow &~\heva{&(x+1)^{2}+ (y-4)^{2}= (x-5) ^{2}+ (y-6)^{2}\\&( x - 5 ) ^{2}+ ( y - 6 ) ^{2}= ( x - 6 ) ^{2}+ ( y - 3 ) ^{2}}\\
\Leftrightarrow &~\heva{&3x+y=11\\&x-3y=-8}\\
\Leftrightarrow & ~\heva{&x=\dfrac{5}{2}\\&y=\dfrac{7}{2}.}
\end{align*}
Tọa độ của vị trí đặt cờ  là $I\left(\dfrac{5}{2}; \dfrac{7}{2}\right)$.\\
Nếu tính từ vị trí ba bạn đứng thì đi về phía đông $2{,}5$ m sau đó đi lên hướng Bắc $3{,}5$ m sẽ đến vị trí cách đều ba bạn nhất.
}
\end{ex}

\begin{ex}%[Dự án Tex hoá đề Tốt nghiệp 2025 + Phát triển 05 đề]%[Đình Nguyên]%[0D8C2-7]
Có $4$ ngăn trong một giá để sách và $7$ quyển sách khác nhau. An xếp hết $7$ quyển sách này vào $4$ ngăn sao cho mỗi ngăn có ít nhất một quyển sách và không có ngăn nào chứa quá $3$ quyển sách.
Các quyển sách được xếp thẳng đứng thành một hàng ngang từ trái sang phải trong mỗi ngăn. Hai cách xếp được gọi là giống nhau nếu thỏa mãn đồng thời hai điều kiện sau:
\begin{itemize}
\item Với từng ngăn, số lượng quyển sách ở ngăn đó là như nhau trong cả hai cách xếp.
\item Với từng ngăn, thứ tự từ trái sang phải của các quyển sách được xếp là như nhau trong cả hai cách xếp.
\end{itemize}
Tính số cách xếp đôi một khác nhau của bạn An.
% \shortans[oly]{$8064$}
\loigiai{
\textbf{Cách 1:} Đếm trực tiếp.\\
Gọi số sách trong $4$ ngăn lần lượt là $n_1, n_2, n_3, n_4$.
Theo đề bài, ta có $n_1 + n_2 + n_3 + n_4 = 7$ với $1 \le n_i \le 3$ ($i=1, 2, 3, 4$).
Có hai trường hợp về số lượng sách trong các ngăn:
\begin{enumerate}[TH 1.]
\item Một ngăn chứa $3$ quyển, một ngăn chứa $2$ quyển và hai ngăn còn lại mỗi ngăn chứa $1$ quyển.
\begin{itemize}
\item Số cách chọn ngăn chứa $3$ quyển sách: $\mathrm{C}_4^1$ cách.
\item Số cách chọn và xếp $3$ trong $7$ quyển sách vào ngăn đó: $\mathrm{A}_7^3$ cách.
\item Số cách chọn ngăn chứa $2$ quyển sách trong $3$ ngăn còn lại: $\mathrm{C}_3^1$ cách.
\item Số cách chọn và xếp $2$ trong $4$ quyển sách còn lại vào ngăn đó: $\mathrm{A}_4^2$ cách.
\item Số cách xếp $2$ quyển sách cuối cùng vào $2$ ngăn còn lại: $2!$ cách.
\end{itemize}
Suy ra, số cách xếp trong trường hợp này là $\mathrm{C}_4^1 \times \mathrm{A}_7^3 \times \mathrm{C}_3^1 \times \mathrm{A}_4^2 \times 2! = 4 \times 210 \times 3 \times 12 \times 2 = 60\,480$ cách.
\item Ba ngăn mỗi ngăn chứa $2$ quyển và một ngăn chứa $1$ quyển.
\begin{itemize}
\item Số cách chọn ngăn chứa $1$ quyển sách: $\mathrm{C}_4^1$ cách.
\item Số cách chọn và xếp $1$ trong $7$ quyển sách vào ngăn đó: $\mathrm{A}_7^1$ cách.
\item Số cách chọn và xếp $2$ trong $6$ quyển sách còn lại vào ngăn thứ nhất trong $3$ ngăn còn lại: $\mathrm{A}_6^2$ cách.
\item Số cách chọn và xếp $2$ trong $4$ quyển sách còn lại vào ngăn thứ hai: $\mathrm{A}_4^2$ cách.
\item Số cách xếp $2$ quyển sách cuối cùng vào ngăn cuối cùng: $\mathrm{A}_2^2$ cách.
\end{itemize}
Suy ra, số cách xếp trong trường hợp này là $\mathrm{C}_4^1 \times \mathrm{A}_7^1 \times \mathrm{A}_6^2 \times \mathrm{A}_4^2 \times \mathrm{A}_2^2 = 4 \times 7 \times 30 \times 12 \times 2 = 20\,160$ cách.
\end{enumerate}
Suy ra tổng số cách xếp là $T=60\,480 + 20\,160 = 80\,640$ cách.\\
Vậy $\dfrac{T}{10}=8\,064$.\\
\textbf{Cách 2:} Ta giải bài toán bằng phương pháp bù trừ qua các bước sau:\\
Đầu tiên, ta tính tổng số cách xếp mà mỗi ngăn có ít nhất 1 quyển sách.\\
Ta sử dụng phương pháp \lq\lq vách ngăn\rq\rq\ cho các vật phẩm phân biệt.
\begin{itemize}
\item Sắp xếp $7$ quyển sách khác nhau thành một hàng ngang: có $7!$ cách.
\item Với $7$ quyển sách xếp thành hàng, ta có $6$ vị trí trống ở giữa. Để chia $7$ quyển sách thành $4$ phần không rỗng (cho $4$ ngăn), ta đặt $3$ vách ngăn vào $3$ trong $6$ vị trí trống đó. Số cách đặt vách ngăn là $\mathrm{C}_6^3$.
\end{itemize}
Tổng số cách xếp sao cho mỗi ngăn có ít nhất một quyển là: $7! \times \mathrm{C}_6^3 = 5040 \times 20 = 100\,800$ cách.\\
Bây giờ ta tính số cách xếp vi phạm điều kiện \lq\lq không có ngăn nào quá 3 quyển\rq\rq.\\
Một cách xếp vi phạm khi có ít nhất một ngăn chứa $4$ quyển trở lên.\\ Vì tổng số sách là $7$ và mỗi ngăn phải có ít nhất $1$ quyển, trường hợp vi phạm duy nhất có thể xảy ra là khi số sách trong các ngăn được chia theo tỉ lệ $(4, 1, 1, 1)$.
\begin{itemize}
\item Chọn $1$ trong $4$ ngăn để chứa $4$ quyển sách: $\mathrm{C}_4^1$ cách.
\item Chọn và xếp $4$ trong $7$ quyển sách vào ngăn đó: $\mathrm{A}_7^4$ cách.
\item Xếp $3$ quyển sách còn lại vào $3$ ngăn còn lại: $3!$ cách.
\end{itemize}
Số cách xếp vi phạm là $\mathrm{C}_4^1 \times \mathrm{A}_7^4 \times 3! = 4 \times 840 \times 6 = 20\,160$ cách.\\
Vậy số cách xếp thỏa mãn các điều kiện của đề bài là
\[T=100\,800 - 20\,160 = 80\,640 \text{ cách.}\]
% Suy ra $\dfrac{T}{10}=8\,064$.
}
\end{ex}

\Closesolutionfile{ans}
\Closesolutionfile{ansbook}
